% Page 034 — Alban pendant la guerre de Cent Ans
% Transcription de la page fournie.
% La pagination de travail suit le classement retenu par l'auteur du projet.

\begin{center}
    {\Large\itshape Alban pendant la guerre de Cent Ans}
\end{center}

\vspace{1.5em}

\noindent
Après Paulin, c'est Curvalle qui, presque à la même date, va être victime à son tour
du brigandage. En effet, en mai 1384, un chef de bande nommé Espagnolet s'empare par
deux fois du château de Curvalle et met en coupe réglée le pays environnant. Les
méfaits sont si grands qu'il faut avoir recours aux forces royales pour se débarrasser
de cet encombrant personnage. En mai 1385, une armée royale commandée par le Seigneur
Gaucher de Passac arrive en Albigeois. Les Seigneurs du voisinage apportent leur
concours pour grossir la troupe. Celle-ci est approvisionnée par Albi. En juin 1385,
elle traverse Alban et va mettre le siège devant le château de Curvalle. Cela ne va
pas tout seul. Ce n'est qu'en octobre que la petite place forte est prise. Tous les
brigands qui s'y sont réfugiés sont pendus haut et court et séance tenante. Pendant
toutes les opérations du siège, les convois de vivres venant d'Albi ont passé par
Alban. On devine que tout ce va et vient n'a pas toujours été de tout repos pour nos
ancêtres Albanais. Le chroniqueur Froissart a raconté très pittoresquement les
péripéties de ce siège. (Les érudits y trouveront plaisir).

\medskip

\noindent
Cette fin de siècle n'est pas de tout repos. Albi s'inquiète et ses consuls viennent
se renseigner auprès de ceux d'Alban. Une fois de plus, des hommes d'armes ont occupé
Pousthomis. Une première démarche a eu lieu en avril 1387. En octobre de la même
année, deux consuls d'Albi arrivent chez nous pour obtenir des renseignements. La
guerre est à nos portes. Ces deux magistrats ont nom Gaudetru et Clergue. Ils vont
se rendre à cheval jusqu'à Nîmes pour y débattre des affaires de leur ville. Ils
sont de retour en nos murs douze jours après. En mai 1388, Clergue renouvelle sa
chevauchée et ne met que dix jours pour ce voyage éclair Albi-Nîmes-Albi.
